\section{ELETTROSTATICA DEI CONDUTTORI}

\subsection{I conduttori}

In fisica le distinzioni non sono mai nette: come, per esempio, non \`e netta la
distinzione fra solido e liquido, cos\`i non lo \`e quella fra conduttore ed
isolante. Siccome la conduzione elettrica \`e dovuta alla mobilit\`a dei
portatori di carica, si pu\`o solo parlare di maggiore o minore mobilit\`a dei
portatori di carica di una data sostanza, ed inoltre questa mobilit\`a pu\`o in
taluni casi dipendere fortemente dalla temperatura, dal trattamento sub\`ito dal
corpo, eccetera. La distinzione che noi amiamo fare in conduttori ed isolanti
\`e quindi una delle tante schematizzazioni che si fanno in fisica, una
schematizzazione, per\`o, che \`e ampiamente giustificata dai numeri: le
mobilit\`a dei portatori di carica in quelli che noi chiamiamo ``buoni
conduttori'' e quella nei ``buoni isolanti'' possono differire di molti, molti,
ordini di grandezza. Questo, in termini concreti, significa che in una scala di
tempi ragionevole, i tempi necessari alle cariche per portarsi in condizioni di
equilibrio possono essere considerati trascurabili nei buoni conduttori, mentre
sono ``infiniti'' nei buoni isolanti.\\
In questo capitolo ci occuperemo solo di conduttori solidi, cio\`e metallici; i
metalli hanno una struttura cristallina costituita dagli ioni degli atomi che la
compongono, ed in questo reticolo gli elettroni, tipicamente uno o due per
atomo, possono muoversi liberamente: un metallo \`e sostanzialmente una grossa
molecola in cui molti atomi uguali si aggregano mettendo a disposizione
dell'intera struttura alcuni dei loro elettroni, cos\`i come in una molecola di
idrogeno gli elettroni non appartengono pi\`u ai singoli nuclei da cui
provengono, ma all'intera molecola. Possiamo quindi immaginare un conduttore
come un gas di elettroni liberi, liberi in quanto la reciproca repulsione
coulombiana \`e in larga misura schermata dalla carica positiva del reticolo in
cui si muovono.\\
Nella maggior parte dei casi i portatori di carica (o \emph{cariche libere}) in
un conduttore sono -- come ci si aspetta -- gli elettroni, ma vedremo che ci
sono eccezioni: nello zinco, per esempio, i portatori di carica sono positivi
(!?), come in molti semiconduttori.

Ci occuperemo -- come dice il titolo di questo capitolo -- dell'elettrostatica
macroscopica dei conduttori: \emph{elettrostatica}, perch\`e ci interessano le
situazioni in cui i portatori di carica sono in quiete nel conduttore, e quindi
in condizioni di equilibrio; \emph{macroscopica}, perch\`e sia all'interno che
all'esterno dei conduttori siamo interessati al campo macroscopico (o medio), e
quindi anche le densit\`a di carica nel conduttore alle quali faremo
riferimento, sono le densit\`a medie (o macroscopiche).\\
La condizione di equilibrio significa che se sui portatori di carica non si
esercitano altre forze, all'interno di un conduttore il campo elettrico deve
essere nullo, altrimenti le cariche libere non sarebbero in quiete. In generale
questo non \`e vero: lasciamo perdere la forza di gravit\`a, oppure la forza
centrifuga che si esercita sugli elettroni di un conduttore in rotazione,
perch\`e queste forze sono insignificanti, ma se il conduttore non \`e omogeneo
(per esempio \`e composto da metalli diversi, come rame e zinco) esistono forze
sugli elettroni che sbrigativamente classifichiamo di origine chimica, ma che in
realt\`a derivano prevalentemente da effetti quantistici. Quindi supporremo che
i conduttori siano omogenei. Oppure il conduttore potrebbe essere in moto
immerso in un campo
magnetico: il campo magnetico esercita sulle cariche in moto una forza (forza di
Lorentz) data da $q(\vec v/c)\wedge\vec B$ che, in condizioni stazionarie, deve
essere compensata dalla presenza di un campo elettrico. Quindi, se non
specifichiamo il contrario, i conduttori saranno supposti in quiete.\\
Per la stessa ragione per cui il campo elettrico deve essere nullo dentro il
conduttore, esso pu\`o essere non nullo sulla superficie, ma deve essere
ortogonale alla superficie: gli elettroni, per i motivi che spiegheremo tra
poco, non possono spontaneamente uscire dal conduttore. Quindi
\begin{itemize}
  \item[i)] un conduttore \`e una regione equipotenziale e la sua superficie \`e
    una superficie equipotenziale;
  \item[ii)] in ogni punto interno al conduttore
    \begin{equation}
      \vec E=0\Rightarrow\ \dive\vec E=0\quad\Rightarrow\quad\rho=0
      \tag{2.1}
    \end{equation}
    quindi in un conduttore, sia esso carico o no, ed eventualmente in presenza
    di cariche esterne, le cariche si dispongono esclusivamente sulla
    superficie;
  \item[iii)] per la (1.87), siccome dentro il conduttore il campo elettrico \`e
    nullo, se $\hat n$ \`e la normale alla superficie del conduttore nel punto
    $P$ orientata verso l'esterno, si ha
    \begin{equation}
      \vec E(P)=E_n(P)\hat n\,;\quad E_n(P)=4\pi\sigma(P)
      \tag{2.2}
    \end{equation}
    quindi le linee di forza sono uscenti dal conduttore nei punti in cui
    $\sigma(P)>0$, entranti dove $\sigma(P)<0$. Anche se la (2.2) lega il campo
    elettrico in un punto alla densit\`a di carica nello stesso punto, il campo
    elettrico \`e determinato da \underline{tutta} la distribuzione di carica e,
    eventualmente, anche dalle cariche esterne.
\end{itemize}
Le linee di forza in presenza di conduttori hanno particolari propriet\`a
derivanti dal fatto che il campo elettrico \`e conservativo. Per esempio non \`e
possibile che fra due conduttori ci siano linee di forza che li connettono in
entrambi i versi, oppure che una stessa linea di forza esca da un conduttore
(dove $\sigma>0$) e ritorni sullo stesso conduttore (in un punto in cui
$\sigma<0$): infatti lungo una linea di forza il potenziale decresce, quindi le
linee di forza vanno sempre da un conduttore a potenziale maggiore ad uno a
potenziale minore. Anche il teorema di Gauss vuole la sua parte: un tubo di
flusso uscente da una porzione di un conduttore sulla quale \`e presente una
carica $q$, se termina su un altro conduttore (o anche su pi\`u conduttori:
basta che non vada, anche in parte, all'infinito), la porzione di conduttore su
cui arriva deve contenere la carica $-q$: basta applicare il teorema di Gauss al
tubo di flusso prolungato sia all'interno del conduttore di partenza che di
quello di arrivo e chiuso da basi interne ai conduttori, dove il campo elettrico
\`e nullo.

\paragraph{Esempi}
\begin{enumerate}
  \item Consideriamo una lastra piana conduttrice infinita, scarica e in assenza
    di sorgenti esterne. A causa della struttura cristallina di un conduttore, a
    livello microscopico la densit\`a di carica sulla superficie ($x=0$) della
    lastra ha una struttura periodica bidimensionale che schematizziamo nel modo
    seguente: $\sigma(y,z)=\sigma_0\cos\kappa y\cos\kappa z$, dove $2\pi/\kappa$
    \`e il periodo, cio\`e il passo reticolare, e quindi dell'ordine di
    $10^{-8}$\,cm (in
    realt\`a un conduttore \`e costituito da microcristalli disposti in modo
    disordinato). Si noti che la media di $\sigma(y,z)$ su regioni di dimensioni
    grandi rispetto al passo reticolare \`e nulla. Vogliamo calcolare il campo
    elettrico all'esterno del conduttore ($x>0$), per vedere l'effetto della
    struttura microscopica della distribuzione di carica.\\
    Data la somiglianza con il problema dell'esempio 1 alla fine del capitolo
    1.10, possiamo sperare che un potenziale della forma
    \[
      \varphi(x,y,z)=\varphi_0e^{-\alpha x}\cos\kappa y\cos\kappa z
      \quad\ (x\ge0)
    \]
    risolva il problema: si ha
    \[
      \left\{\begin{aligned}
        E_x&=\alpha\varphi_0e^{-\alpha x}\cos\kappa y\cos\kappa z\\
        E_y&=\kappa\varphi_0e^{-\alpha x}\sen\kappa y\cos\kappa z\\
        E_z&=\kappa\varphi_0e^{-\alpha x}\cos\kappa y\sen\kappa z
      \end{aligned}\right.
    \]
    e (per $x>0$) \ $\dive\vec E=0$ \ se \ $\alpha=\sqrt2\kappa$, quindi
    \[
      \begin{aligned}
        4\pi\sigma=E_x(0^+)-E_x(0^-)
        &=2E_x(0^+)=2\sqrt2\kappa\varphi_0\cos\kappa y\cos\kappa z\\
        \Rightarrow\quad\varphi_0&=\sqrt2\pi\sigma_0/\kappa
      \end{aligned}
    \]
    Il campo elettrico sulla superficie vale
    $2\pi\sigma_0\cos\kappa y\cos\kappa z$ e siccome
    \[
      \sigma_0\approx e/(2\pi/\kappa)^2
      \simeq4.8\times10^{-10}/10^{-16}=4.8\times10^6\,\mathrm{ues}
    \]
    raggiunge valori molto elevati: tuttavia, come per la densit\`a di carica,
    la sua media \`e nulla ed inoltre, a causa del fattore esponenziale, a
    distanza di pochi passi reticolari si \`e gi\`a ridotto praticamente a zero.
  \item Un disco conduttore di raggio $R$ ruota con velocit\`a angolare costante
    $\omega$ attorno al proprio asse.\\
    Vogliamo calcolare il campo elettrico e la densit\`a di carica dentro il
    conduttore.\\
    In condizioni stazionarie:
    \[
      -e\vec E+m\omega^2\vec r=0\quad\Rightarrow
      \quad\vec E=\frac me\,\omega^2\vec r\,;\qquad4\pi\rho=\dive\vec E
      \quad\Rightarrow\quad\rho=\frac{m\omega^2}{2\pi e}\,\cdot
    \]
    Siccome il disco \`e neutro, deve essere presente anche una densit\`a
    superficiale $\sigma$, che per\`o non \`e possibile calcolare in modo
    elementare. Supponiamo che il disco compia 100 giri/s; calcoliamo il campo
    elettrico alla distanza di 10\,cm dall'asse (la massa dell'elettrone \`e
    $m_e=0.9\times10^{-27}$\,g):
    \[
      E=\frac{0.9\times10^{-27}}{4.8\times10^{-10}}\cdot(2\pi)^2\cdot10^4\cdot10
      \approx10^{-11}\,\mathrm{ues}\,.
    \]
\end{enumerate}
Prima di concludere questo paragrafo, cerchiamo di capire -- almeno
qualitativamente -- perch\`e gli elettroni non scappano dal conduttore.
All'interno del conduttore ogni elettrone, circondato sia da cariche positive
che negative, risente in media di una forza nulla; ci\`o non \`e pi\`u vero
quando si avvicina alla superficie: arrivando dall'interno incontra l'ultimo
strato di cariche positive (gli ioni), circondati da una ``nuvola'' di cariche
negative. Quindi il conduttore ``termina'' con uno strato di carica positiva
segu\`ito da uno di carica negativa: questa distribuzione pu\`o essere
schematizzata come
due piani con cariche superficiali opposte $\sigma$ e $-\sigma$, (detta,
appunto, \emph{doppio strato}) a distanza dell'ordine del passo reticolare $d$,
con il piano negativo verso l'esterno del conduttore, e quindi con una
differenza di potenziale negativa fra il piano pi\`u esterno e quello pi\`u
interno. Siccome gli elettroni sono negativi, il doppio strato costituisce per
essi una barriera di energia potenziale che possiamo stimare dell'ordine di
\[
  \Delta U=e\,\sigma d\approx e\cdot\frac{e}{d^2}\cdot d
  =\frac{e^2}{d}\approx\frac{(4.8\times10^{-10})^2}{10^{-8}}
  \approx2\times10^{-11}\,\mathrm{erg}\approx10\,\mathrm{eV}
\]
(in realt\`a \`e pi\`u piccola per circa un ordine di grandezza).\\
Inoltre, quando un elettrone ha attraversato la superficie, risente da parte del
conduttore di una forza attrattiva dovuta sia al fatto che il conduttore -- con
un elettrone in meno -- \`e carico positivamente, ma soprattutto perch\`e
l'elettrone esterno induce sulla superficie del conduttore prossima
all'elettrone una carica positiva. Complessivamente, l'energia che occorre
fornire all'elettrone per estrarlo dal conduttore \`e detta \emph{lavoro di
estrazione}, ed \`e dell'ordine dell'elettronvolt.\\
L'estrazione degli elettroni si pu\`o realizzare, per esempio, riscaldando il
metallo (\emph{effetto termoionico}), oppure inviando su di esso luce di
lunghezza d'onda opportuna (cio\`e minore di un certo valore di \emph{soglia})
(\emph{effetto fotoelettrico}).\\
Allora si capisce anche che se accostiamo due metalli diversi, che quindi hanno
diversi potenziali di estrazione, i due conduttori non sono allo stesso
potenziale, ma fra essi c'\`e una differenza di potenziale uguale alla
differenza dei potenziali di estrazione, e di conseguenza un forte campo
elettrico nella regione (di qualche $10^{-8}$\,cm) di contatto fra i due. La
citata natura quantistica di questo fatto, anche se apparentemente lo abbiamo
spiegato in termini puramente classici, sta molto a monte: \`e la stessa
struttura cristallina del conduttore, e poi il fatto che nel fenomeno di
aggregazione degli atomi si liberino degli elettroni, che richiede la Meccanica
Quantistica per essere spiegata.\\
Il problema pu\`o essere visto -- formalmente -- come un problema di analisi
dimensionale: supponiamo che le equazioni della Fisica Classica (i principi
della Meccanica pi\`u le equazioni dell'elettromagnetismo) siano in grado di
spiegare che se mettiamo insieme molti atomi, p.es.\ di rame, questi si devono
organizzare in un cristallo con un ben preciso passo reticolare. Quindi le
equazioni devono determinare anche questo passo reticolare in termini delle
costanti che entrano nelle suddette equazioni (massa e carica dell'elettrone,
massa dei nuclei e velocit\`a della luce $c$): l'unica grandezza con le
dimensioni di una lunghezza che pu\`o essere costruita con le costanti
fondamentali \`e il cosiddetto ``raggio classico dell'elettrone''
\ $r_e\equiv e^2/m_ec^2\simeq2.8\times10^{-13}$\,cm, troppo piccolo per poter
essere di qualche utilit\`a per spiegare passi reticolari 5 ordini di grandezza
maggiori. In generale, si pu\`o dire che alla Fisica Classica manca una
lunghezza fondamentale dell'ordine di $10^{-8}$\,cm per spiegare la struttura
della materia, cio\`e le dimensioni degli atomi, delle molecole, i passi
reticolari\dots. In meccanica quantistica interviene un'altra costante
fondamentale, la costante di Planck $\hbar$ che ha le dimensioni di una
quantit\`a di moto per una lunghezza (o di un'energia per un tempo), e con essa
\`e possibile costruire una lunghezza fondamentale, chiamata ``raggio di Bohr''
\ $a_0\equiv\hbar^2/m_ee^2=0.53\times10^{-8}$\,cm, ed \`e proprio la grandezza
che determina dimensioni degli atomi, energie di ionizzazione, lavori di
estrazione, eccetera.

\subsection{Il fenomeno dell'induzione ed il problema dell'elettrostatica dei
conduttori}

Se avviciniamo una carica ad un conduttore scarico, sulla sua superficie si
manifesta una distribuzione di cariche: di segno opposto in prossimit\`a della
carica esterna, e quindi (il conduttore \`e scarico) dello stesso segno pi\`u o
meno dalla parte opposta. Questo \`e il fenomeno dell'\emph{induzione
elettrostatica}; intuitivamente si capisce che esso \`e dovuto alle forze
coulombiane che la carica esterna esercita sulle cariche libere nel conduttore,
ma non dimentichiamoci che esistono anche le forze coulombiane fra le cariche
libere. Per questo motivo \`e pi\`u suggestivo, e soprattutto pi\`u utile,
vedere la cosa in termini diversi: le cariche si devono ridistribuire sulla
superficie del conduttore in modo tale che la somma del campo elettrico da esse
prodotto e quello prodotto dalla carica esterna sia nullo all'interno del
conduttore, e quindi il conduttore sia equipotenziale.\\
Questo ci fa capire che il problema dell'elettrostatica in presenza di
conduttori \`e tutt'altro che banale: si pensi ad esempio ad una carica
puntiforme in prossimit\`a di una sfera conduttrice; come si fa a sapere quale
deve essere la distribuzione superficiale sulla sfera in grado di garantire
campo nullo dentro tutta la sfera? (Pi\`u avanti saremo in grado di risolvere
questo problema).

Un problema che invece siamo in grado di risolvere subito \`e quello di
determinare la distribuzione di carica indotta sulla superficie di una lastra
conduttrice piana ``infinita'' da una carica puntiforme $q$: il campo \`e
ovunque la sovrapposizione del campo della carica $q$ e di quello della
distribuzione (non uniforme) incognita $\sigma$ sulla superficie della lastra
affacciata alla carica $q$; supponiamo (ma pi\`u avanti vedremo che questa
ipotesi non \`e necessaria) che lo spessore della lastra sia infinito (per la
precisione, ``pi\`u infinito'' della larghezza della lastra), per cui se
eventualmente ci fossero delle cariche distribuite anche sulla superficie
``lontana'' queste non contribuirebbero. Dentro la lastra il campo \`e nullo e
quindi quello prodotto dalla distribuzione $\sigma$ \`e opposto a quello della
carica $q$ (figura 2.1); nell'attraversare la distribuzione cambia segno la
componente normale alla superficie del campo prodotto da $\sigma$ (invarianza
per riflessioni), mentre resta inalterata la componente parallela, quindi il
campo prodotto da $\sigma$ sulla superficie esterna della lastra \`e uguale a
quello di una carica $-q$, immagine speculare di $q$ rispetto alla superficie
della lastra: noto il campo sulla superficie della lastra (che \`e due volte la
componente normale di quello dovuto alla carica $q$), per la (2.2) \`e nota la
densit\`a $\sigma$.
\fig{fig2_1}{2.1}
Vedremo pi\`u avanti che in tutti i punti esterni alla lastra (dalla parte di
$q$), e non solo sulla superficie della lastra, il campo \`e uguale a quello
delle due cariche $q$ e $-q$: questo, come pure quello relativo allo spessore
della lastra, sono due problemi che finch\'e non abbiamo le idee chiare su come
si pone il problema dell'elettrostatica in presenza di conduttori, dobbiamo
lasciare in sospeso.

Vediamo allora come si formula il problema dell'elettrostatica in presenza di
conduttori, cio\`e quali sono i dati e quali le incognite. Contrariamente al
caso delle cariche fisse, qui noi non abbiamo nessun controllo sulla
distribuzione delle cariche sui conduttori:
ci\`o che noi possiamo decidere \`e solo quant'\`e il potenziale di un
conduttore, in quanto con opportune batterie possiamo portarlo al potenziale che
ci pare; oppure la carica che gli cediamo, poi come questa si distribuisce non
dipende da noi. Quindi, dati $n$ conduttori di forma nota e in posizioni note
(figura 2.2), il problema dell'elettrostatica dei conduttori si pu\`o presentare
nelle seguenti forme:
\begin{itemize}
  \item[i)] \textbf{Problema di Dirichlet:} sono dati i potenziali
    $\varphi_i$\,, $i=1\cdots n$ \ di tutti i conduttori. Il problema \`e quello
    di determinare il campo elettrico (o il potenziale) in tutto lo spazio, la
    carica totale su ciascun conduttore e come \`e distribuita.\\
    In effetti l'unico problema \`e quello della determinazione del potenziale,
    in quanto noto il potenziale \`e noto il campo elettrico, quindi -- per la
    (2.2) -- le distribuzioni $\sigma_i(x,y,z)$ su tutti i conduttori e quindi
    le loro cariche $Q_i=\int_{S_i}\sigma_i(x,y,z)\,dS_i$.
  \item[ii)] \textbf{Problema di Neumann:} sono date le cariche $Q_i$\,,
    $i=1\cdots n$ \ di tutti i conduttori. Il problema \`e sempre quello di
    determinare il potenziale in tutto lo spazio, sapendo che assume valori
    costanti (incogniti) sui conduttori.
  \item[iii)] \textbf{Problema misto:} di alcuni conduttori sono noti i
    potenziali, di altri le cariche.
\end{itemize}
Vedremo che sia il problema di Neumann che quello misto sono riconducibili al
problema di Dirichlet, per cui cominciamo con questo.\\
Siano dati $n$ conduttori $c_i$, siano $S_i$ le loro superfici esterne e $S_0$
una superficie che li racchiude (figura 2.2), che pu\`o essere la superficie
interna di un conduttore cavo, oppure la superficie all'infinito (che possiamo
pensare come la superficie sferica interna di un grande conduttore cavo). Se non
esistono cariche (fisse) esterne ai conduttori, nella regione $\mathcal R$ di
spazio delimitata dalle superfici $S_i$ dei conduttori e da $S_0$ il potenziale
deve soddisfare l'equazione di Laplace \ $\Delta\,\varphi(x,y,z)=0$ \ con
assegnate condizioni al contorno: \ $\varphi(x,y,z)=\varphi_i$ \ sull'i--esimo
conduttore, \ $\varphi(x,y,z)=\varphi_0$ \ sulla superficie $S_0$; di solito,
dato che il potenziale \`e definito a meno di una costante, $\varphi_0$ viene
preso uguale a zero, specialmente se $S_0$ \`e la superficie all'infinito:
sovente $S_0$ viene indicata come la \textbf{terra} o la \textbf{massa} del
sistema; ``mettere a terra'' (o ``a massa'') un conduttore significa portarlo al
potenziale di $S_0$ stabilendo tra i due il contatto tramite un filo conduttore.
\fig{fig2_2}{2.2}
Date queste condizioni al contorno, \`e noto che il problema ammette soluzione:
non possiamo dimostrarlo, ma sappiamo che allora la soluzione \`e unica.\\
Se ci sono anche delle cariche fisse, per fissare le idee diciamo una
distribuzione $\rho$, il potenziale deve soddisfare (in $\mathcal R$)
l'equazione di Poisson \ $\Delta\,\varphi(x,y,z)=-4\pi\rho(x,y,z)$ \ con le
condizioni al contorno $\varphi(x,y,z)=\varphi_i$ \quad $(i=0\cdots n)$:
mediante la (1.59) calcoliamo il potenziale $\tilde\varphi(x,y,z)$ dovuto solo
alle cariche fisse, e sia $\tilde\varphi_i(x,y,z)$ il valore che assume su $S_i$
\quad $(i=0\cdots n)$. Quindi risolviamo l'equazione di Laplace
$\Delta\,\chi(x,y,z)=0$ in $\mathcal R$, con le condizioni al contorno sulle
superfici $S_i$ \quad $\chi_i(x,y,z)=\varphi_i-\tilde\varphi_i(x,y,z)$: anche in
questo caso la soluzione esiste ed \`e unica, e
\ $\varphi(x,y,z)=\chi(x,y,z)+\tilde\varphi(x,y,z)$ \ risolve il problema.

Quindi il problema di Dirichlet \`e un problema prettamente matematico: la
soluzione
sappiamo che esiste, ma trovarla \`e tutto un altro discorso; non esiste una
formula risolutiva (come la (1.59) per il caso delle cariche fisse), per cui i
problemi risolubili sono pochi, e acquistano particolare importanza metodi di
soluzione applicabili in casi particolari (come il metodo delle immagini), e
anche metodi generali di approssimazione (come il metodo del rilassamento e il
calcolo delle variazioni).

Come gi\`a nel caso delle cariche fisse, le chances di risolvere un problema
aumentano moltissimo in presenza di simmetrie (conduttori sferici, cilindrici,
piani conduttori infiniti\dots); per\`o \`e chiaro che non possiamo fare
riferimento alla discussione del capitolo 1.6, in quanto presuppone la
conoscenza della distribuzione di carica: dobbiamo riformulare il teorema
relativo alle leggi di trasformazione del campo elettrico ed il corollario sul
campo delle distribuzioni simmetriche in termini del potenziale (anzich\`e di
campo elettrico) e delle condizioni al contorno (anzich\`e della distribuzione
di carica). Per semplicit\`a enunciamo il teorema nell'ipotesi che non esistano
cariche fisse esterne ai conduttori: la formulazione generale \`e possibile ed
\`e immediata.

\textbf{\textit{Teorema.}} \textit{Sia $\varphi(x)$
\ $\big(x=(x_1,x_2,x_3)\big)$ la soluzione del problema di Dirichlet nella
regione $\mathcal R$ delimitata dalle superfici $S_i$ dei conduttori a
potenziali $\varphi_i$, e dalla superficie $S_0$ a potenziale $\varphi_0$, e sia
$\psi(x)=Rx+a$ una isometria dello spazio \big(v.\ 1.37)\big). Allora
$\varphi'(x)=\varphi\big(\psi^{-1}(x)\big)$ \`e la soluzione del problema di
Dirichlet nella regione $\mathcal R'=\psi(\mathcal R)$ delimitata dalle
superfici $S_i'=\psi(S_i)$ di conduttori a potenziali $\varphi_i$, e dalla
superficie $S_0'=\psi(S_0)$ a potenziale $\varphi_0$. (Se $S_0$ \`e la
superficie all'infinito, $S_0'\equiv S_0$).}

Segue allora dal teorema, ma solo a posteriori, che
$\sigma_i'(x)=\sigma_i\big(\psi^{-1}(x)\big)$.

Basta dimostrare che $\varphi'(x)$ \`e una funzione armonica nella regione
$\mathcal R'$, dopodich\`e il teorema \`e dimostrato in quanto, siccome
$\varphi'\big(\psi(x)\big)=\varphi(x)$, \`e ovvio che $\varphi'(x)$ assume i
valori $\varphi_i$ su $S_i'$ e $\varphi_0$ su $S_0'$, quindi \`e una soluzione
del problema di Dirichlet in $\mathcal R'$, e per l'unicit\`a della soluzione
\`e \underline{la} soluzione. Per quanto riguarda la dimostrazione che se
$\Delta\,\varphi(x)=0$, allora $\Delta\,\varphi'(x)=0$, si pu\`o fare in modo
elementare applicando le regole di derivazione (parziale) delle funzioni
composte. Ma \`e meno laborioso e pi\`u istruttivo ragionare come segue:
$\Delta\,\varphi'(x)$ \`e, per la (1.83), il flusso per unit\`a di volume del
campo vettoriale $\grad\varphi'(x)$ nel punto $x\in\mathcal R'$, ma questo \`e
uguale al flusso per unit\`a di volume (i volumi sono invarianti per isometrie)
del campo vettoriale $\grad\varphi(x)$ nel punto $\psi^{-1}(x)\in\mathcal R$,
che \`e nullo per ipotesi, quindi \`e nullo anche il primo.\\
Segue allora immediatamente il

\textbf{\textit{Corollario.}} \textit{Se $S_i'=S_i$, \ $i=0\cdots n$\,, \ allora
$\varphi(x)=\varphi\big(\psi^{-1}(x)\big)$.}

Cio\`e: se una isometria lascia invariate le superfici equipotenziali $S_i$,
anche il potenziale resta invariato. (Per maggiore generalit\`a $S_i$ \`e
l'insieme delle superfici su cui $\varphi=\varphi_i$: possono esistere pi\`u
conduttori identici allo stesso potenziale, mandati l'uno nell'altro dalla
trasformazione).

\paragraph{Esempio}
\begin{quote}
  Sfera conduttrice di raggio $R$ a potenziale $\overline\varphi$. Vogliamo
  determinare il potenziale in tutto lo spazio esterno alla sfera.\\
  Possiamo ragionare in due modi: la funzione
  \begin{equation}
    \varphi(x,y,z)=\overline\varphi\,\frac Rr
    \tag{2.3}
  \end{equation}
  soddisfa l'equazione di Laplace per $r>R$ e la corretta condizione al
  contorno, quindi -- per l'unicit\`a della soluzione -- \`e la soluzione.
  Questo metodo funziona quando \`e facile indovinare la soluzione.\\
  Oppure: siccome il conduttore \`e invariante per rotazioni, per il corollario
  tale deve essere anche il potenziale, quindi il campo elettrico \`e un campo
  centrale e di conseguenza la densit\`a di carica $\sigma$ \`e costante sulla
  superficie della sfera. Allora per la (1.69)
  \begin{equation}
    \varphi(r)=\frac{4\pi\sigma R^2}{r}
    \quad\ r\ge R\qquad\quad\varphi(R)=\overline\varphi
    \ \Rightarrow\ Q=4\pi\sigma R^2=\overline\varphi R
    \tag{2.4}
  \end{equation}
  in accordo con la (2.3).
\end{quote}

\subsection{Le gabbie di Faraday}

In questo paragrafo utilizzeremo l'equivalenza fra il problema di Dirichlet, il
problema di Neumann e quello misto, equivalenza che dimostreremo nel prossimo
paragrafo.\\
Sia dato un conduttore cavo, sia $S_\mathrm{i}$ la sua superficie interna ed
$S_\mathrm{e}$ quella esterna. Consideriamo diversi casi.
\begin{enumerate}
  \item \emph{All'interno della cavit\`a non sono presenti cariche.}\\
    Il potenziale \`e costante ($\varphi_0$) sul conduttore, e quindi su
    $S_\mathrm{i}$, e soddisfa l'equazione di Laplace in tutti i punti interni:
    $\varphi(x,y,z)=\varphi_0$ soddisfa l'equazione e la condizione al contorno
    e quindi \`e la soluzione. Quindi il campo elettrico \`e nullo dentro la
    cavit\`a e, per la (2.2), non ci sono cariche sulla superficie interna,
    indipendentemente dal fatto che il conduttore sia carico o no, e dal fatto
    che all'esterno ci sia o meno un campo elettrico. Se il conduttore \`e
    carico, tutta la carica si trova sulla superficie esterna.\\
    Se la sfera dell'esempio alla fine del paragrafo precedente \`e cava,
    qualunque sia la forma della sua superficie interna, la carica della sfera
    \`e solo quella sulla superficie esterna data dalla (2.4), ed il potenziale
    dentro la cavit\`a \`e $\overline\varphi$.
  \item \emph{All'interno della cavit\`a sono presenti delle cariche: il
    problema interno.}\\
    Se $Q_\mathrm{i}$ \`e la carica totale contenuta dentro la cavit\`a, questa
    induce su $S_\mathrm{i}$ una carica uguale ed opposta: infatti presa come
    superficie di Gauss una qualsiasi superficie $S$ (figura 2.3) compresa fra
    $S_\mathrm{i}$ ed $S_\mathrm{e}$, e quindi tutta interna al conduttore,
    siccome $\vec E=0$ dentro il conduttore il flusso attraverso $S$ \`e nullo,
    e quindi la carica totale racchiusa da $S$, $Q_\mathrm{i}$ pi\`u la carica
    su $S_\mathrm{i}$, \`e nulla. Se il conduttore \`e scarico, una carica pari
    a $Q_\mathrm{i}$ si ritrova distribuita sulla superficie esterna; se il
    conduttore ha una carica $Q$, su $S_\mathrm{e}$ \`e presente la carica
    $Q^{tot}=Q+Q_\mathrm{i}$. Quindi sulla superficie esterna \`e sempre
    presente tutta la carica contenuta dal conduttore.
    \fig{fig2_3}{2.3}
    Anche in questo caso il potenziale \`e completamente determinato dalle
    condizioni assegnate all'interno (posizione delle cariche fisse, potenziali
    o cariche degli eventuali conduttori) e dal valore (costante) che il
    potenziale assume su $S_\mathrm{i}$. Il potenziale del conduttore, cio\`e
    della ``gabbia di Faraday'' \underline{dipende} dalle condizioni esterne (se
    si richiede
    che $\varphi(\infty)=0$), per cui se cambiano le condizioni esterne il
    potenziale all'interno cambia per una costante (additiva), ma questo \`e
    irrilevante perch\`e ci\`o che conta sono sempre le differenze di
    potenziale, ovvero il campo elettrico. Quindi all'interno della gabbia campo
    elettrico, differenze di potenziale e distribuzioni di carica sugli
    eventuali conduttori sono indipendenti dalle condizioni esterne e dalla
    carica del conduttore. Quindi, in particolare, siccome all'esterno il campo
    potrebbe anche essere nullo, le cariche sulla superficie interna
    $S_\mathrm{i}$ si distribuiscono in modo tale da garantire che il campo
    prodotto da esse e dalle cariche interne \`e nullo nella regione occupata
    dal conduttore.
  \item \emph{All'interno della cavit\`a sono presenti delle cariche: il
    problema esterno.}\\
    La carica totale sulla superficie esterna del conduttore dipende (anche)
    dalla carica interna alla cavit\`a ($Q^\mathrm{tot}=Q+Q_\mathrm{i}$), ma la
    sua distribuzione $\sigma_\mathrm{e}$ su $S_\mathrm{e}$ non dipende da come
    la carica interna $Q_\mathrm{i}$ \`e distribuita all'interno, cio\`e dalle
    posizioni delle cariche fisse e dalle distribuzioni sugli eventuali
    conduttori. Infatti, siccome la carica totale di un conduttore ne determina
    il potenziale (equivalenza del problema di Neumann e di Dirichlet), la
    soluzione all'esterno del conduttore, e quindi anche la distribuzione
    $\sigma_\mathrm{e}$, \`e determinata, e unica, indipendentemente da come le
    cariche sono distribuite all'interno. In particolare, le cariche interne
    possono venire rimosse, ed una carica uguale fornita al conduttore (in modo
    da mantenere inalterata la carica totale), e all'esterno non cambia nulla.
    Quindi anche le cariche sulla superficie esterna insieme alle cariche
    esterne generano un campo nullo nella regione occupata dal conduttore: ci\`o
    significa che il problema interno e quello esterno sono indipendenti.
\end{enumerate}

\paragraph{Esempio}
\begin{quote}
  Una carica puntiforme $q$ \`e posta al centro di una corona sferica
  conduttrice di raggi $R_\mathrm{i}$ e $R_\mathrm{e}$.\\
  L'interno: per ``ovvie ragioni di simmetria'' il campo all'interno \`e
  centrale; per il teorema di Gauss \`e il campo della sola carica $q$. Sulla
  superficie interna $\sigma_\mathrm{i}$ \`e costante e
  $4\pi R_\mathrm{i}^2\sigma_\mathrm{i}=-q$, oppure dalla (2.2)
  $4\pi\sigma_\mathrm{i}=-q/R_\mathrm{i}^2$. Si noti che $\sigma_\mathrm{i}$,
  per la (1.47), produce campo nullo per $r<R_\mathrm{i}$ (e quindi, in questo
  caso, il campo interno \`e, come sapevamo, solo quello generato da $q$), e che
  $\sigma_\mathrm{i}$ insieme a $q$ generano campo nullo per $r>R_\mathrm{i}$.\\
  L'esterno: ci si riconduce facilmente al problema gi\`a discusso di una sfera
  conduttrice con potenziale assegnato.\\
  In definitiva:
  \begin{equation}
    \vec E=\begin{cases}\dfrac{q}{r^2}\,\hat r&0<r<R_\mathrm{i}\\[6pt]
    0&R_\mathrm{i}<r<R_\mathrm{e}\\[2pt]
    \dfrac{q}{r^2}\,\hat r&r>R_\mathrm{e}\end{cases}
    \qquad\quad
    \varphi(r)=\begin{cases}\dfrac qr&r>R_\mathrm{e}\\[6pt]
    \dfrac{q}{R_\mathrm{e}}&R_\mathrm{i}<r<R_\mathrm{e}\\[6pt]
    \dfrac qr+\dfrac{q}{R_\mathrm{e}}-\dfrac{q}{R_\mathrm{i}}
    &0<r<R_\mathrm{i}\,.\end{cases}
    \tag{2.5}
  \end{equation}
  Si noti che avendo richiesto $\varphi(\infty)=0$, la determinazione del
  potenziale ci \`e convenuto farla dall'esterno verso l'interno.\\
  Se la carica $q$ \`e spostata rispetto al centro, cambia il campo all'interno
  (che non \`e pi\`u solo quello generato dalla carica $q$) e la distribuzione
  $\sigma_\mathrm{i}$ non \`e pi\`u uniforme e contribuisce a determinare il
  campo all'interno; campo e distribuzione $\sigma_\mathrm{i}$ non
  possono essere calcolati in modo semplice (per\`o impareremo a farlo), mentre
  all'esterno non cambia nulla.
\end{quote}

\subsection{I coefficienti di capacit\`a e di potenziale}

Sono dati $n$ conduttori $c_i$ nella regione $\mathcal R$, racchiusi dalla
superficie $S_0$: per dimostrare che il problema di Neumann \`e equivalente al
problema di Dirichlet basta stabilire che i potenziali $\varphi_i$ dei
conduttori e le loro cariche $Q_i$ sono legati da una relazione lineare e
invertibile. Poniamo uguale a zero il potenziale di $S_0$.\\
Sia $\varphi'(x,y,z)$ la soluzione del problema di Dirichlet con le condizioni
al contorno $\varphi'(x,y,z)=\varphi_i'$ \ su $c_i$ e siano $Q_i'$ le cariche
dei vari conduttori:
\begin{equation}
  \begin{split}
    Q_i'&=\int_{S_i}\sigma_i'(x,y,z)\,dS_i
    =\frac{1}{4\pi}\int_{S_i}E_n'(x,y,z)\,dS_i\\
    &=-\frac{1}{4\pi}\int_{S_i}\grad\varphi'(x,y,z)\cdot d\vec S_i
  \end{split}
  \tag{2.6}
\end{equation}
che quindi dipendono linearmente da $\varphi'(x,y,z)$.\\
Sia ora $\varphi''(x,y,z)$ la soluzione del problema di Dirichlet con le
condizioni al contorno $\varphi''(x,y,z)=\varphi_i''$ \ su $c_i$ e siano $Q_i''$
le cariche dei conduttori: siccome l'equazione di Laplace \`e lineare ed
omogenea,
\[
  \varphi(x,y,z)=\alpha\varphi'(x,y,z)+\beta\varphi''(x,y,z)
\]
\`e anch'essa soluzione dell'equazione di Laplace, con le condizioni al contorno
\ $\varphi_i=\alpha\varphi_i'+\beta\varphi_i''$, e per la (2.6) le cariche dei
conduttori sono $Q_i=\alpha Q_i'+\beta Q_i''$.\\
Quindi l'applicazione $\{\varphi_i\}\to\{Q_i\}$ \ \`e una applicazione lineare:
\begin{equation}
  Q_i\,=\sum_{j=1}^nC_{ij}\,\varphi_j\,;\qquad\qquad i=1\cdots n.
  \tag{2.7}
\end{equation}
I coefficienti $C_{ij}$ nella (2.7) sono detti \textbf{coefficienti di
capacit\`a} del sistema: ciascuno di essi dipende (solo) dalla forma e dalla
disposizione di \underline{tutti} i conduttori e quindi costituiscono una
caratteristica esclusivamente geometrica (non elettrica) del sistema.\\
La matrice $C$ \`e invertibile: infatti, per la (1.92), la (2.7) e la (1.94) si
ha
\begin{equation}
  U=\frac12\sum_{i=1}^n\varphi_iQ_i=\frac12\sum_{i,j}\varphi_iC_{ij}\varphi_j
  =\frac{1}{8\pi}\int\vec E^2(\vec r\,)\,dV>0
  \tag{2.8}
\end{equation}
cio\`e la forma quadratica associata alla matrice $C$ \`e definita positiva.
Quindi, posto $P\equiv C^{-1}$, si ha
\begin{equation}
  \varphi_i\,=\sum_{j=1}^nP_{ij}\,Q_j\,;\qquad\qquad i=1\cdots n.
  \tag{2.9}
\end{equation}
I coefficienti $P_{ij}$ nella (2.9) sono detti \textbf{coefficienti di
potenziale} del sistema: come i coefficienti di capacit\`a (da cui dipendono)
sono una caratteristica esclusivamente geometrica del sistema.\\
La (2.9) mostra che il problema di Neumann \`e sempre riconducibile al problema
del
Dirichlet, e la stessa cosa vale per il problema misto, in quanto la (2.7) pu\`o
essere risolta rispetto ai soli potenziali incogniti.

Similmente alla (2.8) abbiamo
\begin{equation}
  U=\frac12\sum_{i=1}^nQ_i\varphi_i=\frac12\sum_{i,j}Q_iP_{ij}Q_j.
  \tag{2.10}
\end{equation}
Vediamo alcune delle propriet\`a dei coefficienti $C_{ij}$ e $P_{ij}$:
\begin{itemize}
  \item[i)] La matrice $C$, e quindi anche $P$ \`e simmetrica:
    $C_{ij}=C_{ji}$;\quad $P_{ij}=P_{ji}$.\\
    La dimostrazione si basa sul fatto che l'energia elettrostatica di una
    configurazione di cariche \`e indipendente dal modo in cui la configurazione
    viene realizzata. Per semplicit\`a consideriamo il caso di due soli
    conduttori a potenziali $\varphi_1$ e $\varphi_2$. Riscriviamo per questo
    caso la (2.7):
    \begin{equation}
      \begin{aligned}
        Q_1&=C_{11}\varphi_1+C_{12}\varphi_2\\
        Q_2&=C_{21}\varphi_1+C_{22}\varphi_2\,.
      \end{aligned}
      \tag{2.11}
    \end{equation}
    Per realizzare questa configurazione procediamo nel modo seguente: mettiamo
    a massa il conduttore $c_2$ (cio\`e a potenziale zero), e portiamo $c_1$ al
    potenziale $\varphi_1$ (figura 2.4a); dette $Q_1'$ e $Q_2'$ le cariche sui
    due conduttori, per la (2.11) si ha \ $Q_1'=C_{11}\,\varphi_1$;
    \ $Q_2'=C_{21}\,\varphi_1$. Si noti che, per il fenomeno dell'induzione
    elettrostatica, il conduttore $c_2$, anche se a potenziale zero, ha una
    carica non nulla. Per la (2.8) l'energia $U_1$ della configurazione \`e
    \ $U_1=\frac12C_{11}\varphi_1^2$ \ ed \`e il lavoro fatto per prelevare
    dalla massa ($S_0$) la carica $Q_1'$ e portarla su $c_1$. Non si compie
    lavoro per portare $Q_2'$ su $c_2$ perch\'e $c_2$ \`e allo stesso potenziale
    della massa. Ora teniamo $c_1$ al potenziale $\varphi_1$ (p.es.\ con una
    batteria: figura 2.4b) e portiamo $c_2$ al potenziale $\varphi_2$: sempre
    per la (2.11) al conduttore $c_1$ viene trasferita una carica
    \ $Q_1''=C_{12}\varphi_2$ \ compiendo il lavoro
    \ $Q_1''\varphi_1=\varphi_1C_{12}\varphi_2$ \ (senza il fattore $1/2$: il
    potenziale \`e costantemente $\varphi_1$), mentre al conduttore $c_2$
    occorre trasferire la carica \ $Q_2''=C_{22}\,\varphi_2$ \ compiendo il
    lavoro \ $\frac12Q_2''\varphi_2=\frac12C_{22}\varphi_2^2$. Mettendo tutto
    insieme, il lavoro fatto \`e
    \begin{center}
      \begin{tabular}{cc}
        \includegraphics[scale=1.2]{fig2_4a} &
        \includegraphics[scale=1.2]{fig2_4b}\\
        {\small fig.~2.4a} & {\small fig.~2.4b}
      \end{tabular}
    \end{center}
    \begin{equation}
      U=\frac12C_{11}\varphi_1^2
      +\varphi_1C_{12}\varphi_2+\frac12C_{22}\varphi_2^2
      \tag{2.12}
    \end{equation}
    e confrontando con la (2.8) si ottiene $C_{12}=C_{21}$, e quindi anche
    \ $P_{12}=P_{21}$.\\
    La simmetria dei coefficienti di capacit\`a esprime un fatto tutt'altro che
    ovvio: per esempio nel caso di due conduttori $C_{12}$ \`e la carica del
    conduttore $c_1$ quando il suo potenziale \`e nullo (conduttore a massa) e
    quello del conduttore $c_2$ \`e $\varphi_2=1$, ed \`e uguale alla carica
    $C_{21}$ del conduttore $c_2$ se invece quest'ultimo \`e a potenziale zero e
    $\varphi_1=1$, una situazione tutt'altro che simmetrica della precedente,
    data l'arbitrariet\`a delle posizioni e delle forme dei due conduttori.
  \item[ii)] \begin{equation}
      C_{ii}>0,\ P_{ii}>0;\quad\ C_{ij}<0,\ P_{ij}>0\quad\ (i\neq j).
      \tag{2.13}
    \end{equation}
    La 1$^{\rm a}$ segue immediatamente dalla positivit\`a della matrice $C$; la
    2$^{\rm a}$ dal fatto che anche $P\equiv C^{-1}$ \`e positiva. Per
    dimostrare che $C_{ij}<0$, notiamo che $C_{ij}$, per la (2.7), \`e la carica
    del conduttore $c_i$ quando il conduttore $c_j$ \`e a potenziale
    $\varphi_j=1$, e tutti gli altri conduttori (compreso $c_i$) sono a massa.
    Siccome $C_{jj}>0$, la carica di $c_j$ \`e positiva, e per il fenomeno
    dell'induzione sugli altri conduttori viene indotta una carica negativa.
    Quindi $C_{ij}<0$. Oppure: sia $\varphi_j>0$ e tutti gli altri conduttori
    siano a terra. Siccome $\varphi(x,y,z)$ pu\`o essere minimo solo sui
    conduttori, $\varphi(x,y,z)\ge0$ ed \`e minimo su tutti i conduttori a
    terra, sui quali quindi il campo elettrico \`e ovunque diretto verso
    l'interno, da cui $Q_i<0$\quad$i\neq j$.\\
    In modo analogo si dimostra la positivit\`a dei $P_{ij}$.
\end{itemize}

\paragraph{Esempio}
\begin{quote}
  Una sfera conduttrice (conduttore $c_1$) di raggio $a$ al centro di una corona
  sferica conduttrice (conduttore $c_2$), raggio interno $b$ e raggio esterno
  $c$. Calcoliamo i coefficienti di capacit\`a e di potenziale del sistema.\\
  Cominciamo con i coefficienti di potenziale, perch\'e \`e pi\`u semplice:
  infatti, se $Q_1$ \`e la carica sulla sfera interna e $Q_2$ quella sulla
  corona, dalle (2.5) si ha
  \fig{fig2_5}{2.5}
  \begin{equation}
    \varphi(r)=\begin{cases}\dfrac{Q_1+Q_2}{r}&\text{per }r\ge c\\[10pt]
    \dfrac{Q_1}{r}-\dfrac{Q_1}{b}+\dfrac{Q_1+Q_2}{c}&\text{per }a\le r\le
    b\end{cases}
    \tag{2.14}
  \end{equation}
  da cui
  \begin{equation}
    \begin{gathered}
      \varphi_1=\Big(\frac1a+\frac1c-\frac1b\Big)Q_1+\frac{Q_2}{c}\,;
      \quad\ \varphi_2=\frac{Q_1+Q_2}{c}\qquad\Rightarrow\\
      P_{11}=\frac1a+\frac1c-\frac1b\,;
      \quad\ P_{12}=P_{21}=P_{22}=\frac1c\,\cdot
    \end{gathered}
    \tag{2.15}
  \end{equation}
  Anche i coefficienti di capacit\`a si possono calcolare direttamente, ma a
  questo punto \`e pi\`u facile invertire la matrice $P_{ij}$ e si ottiene
  \begin{equation}
    C_{11}=\frac{ab}{b-a}\,;
    \quad\ C_{21}=C_{12}=\frac{ab}{a-b}\,;\quad\ C_{22}=c+\frac{ab}{b-a}\,\cdot
    \tag{2.16}
  \end{equation}
  Vale la pena di osservare che $C_{11}$ e $C_{22}$ (stesso discorso per
  $P_{11}$ e $P_{22}$) non dipendono solo dalle caratteristiche geometriche del
  conduttore al quale si riferiscono, ma da quelle di tutto il sistema.\\
  Un'altra cosa da notare sono le uguaglianze (una conseguenza dell'altra)
  \ $C_{12}=-C_{11}$, \ $P_{21}=P_{22}$: come vedremo nel prossimo paragrafo
  esse valgono in generale tutte le volte che un conduttore ($c_2$) costituisce
  una gabbia di Faraday per l'altro.
\end{quote}

\subsection{Capacit\`a e condensatori}

Cominciamo con il caso in cui il sistema di conduttori \ $c_1\cdots c_n$ \ \`e
costituito da un solo conduttore: in questo caso la (2.7) si riduce a
\begin{equation}
  Q=C\,\varphi\quad\Rightarrow\quad C=\frac Q\varphi
  \tag{2.17}
\end{equation}
e $C$, il rapporto fra carica del conduttore ed il suo potenziale rispetto alla
terra, \`e detta la \textbf{capacit\`a} del conduttore. Nel sistema di Gauss la
capacit\`a si misura in centimetri: 1\,cm \`e la capacit\`a di un conduttore al
quale occorre fornire la carica di 1 franklin perch\`e il suo potenziale
(rispetto alla terra) sia 1 statvolt. Nel SI l'unit\`a di capacit\`a \`e il
farad (F):
\begin{equation}
  1\,\mathrm F=1\,\mathrm C/1\,\mathrm V=(3\times10^9\cdot300)\,\mathrm{ues}
  =9\times10^{11}\,\mathrm{ues}\approx10^{12}\,\mathrm{ues}\,.
  \tag{2.18}
\end{equation}
Consideriamo ad esempio un conduttore sferico di raggio $a$. La massa
($\varphi_0=0$) sia una sfera cava conduttrice di raggio interno $R$ concentrica
con la sfera.\\
Siccome $\varphi(r)=Q_1/r-Q_1/R$ \quad $(a\le r\le R)$, si ha
\begin{equation}
  \varphi=Q\Big(\frac1a-\frac1R\Big)\quad\Rightarrow\quad C=\frac{aR}{R-a}
  \tag{2.19}
\end{equation}
$C$ \ \`e la capacit\`a della sfera rispetto alla massa (cio\`e la superficie
$S_0$). Se $R\to\infty$, $C\to a$: quindi nel sistema di Gauss 1\,ues \`e la
capacit\`a rispetto all'infinito di una sfera conduttrice di raggio 1\,cm. Nel
SI i coefficienti di capacit\`a sono tutti moltiplicati per $4\pi\epsilon_0$,
per cui la capacit\`a della sfera \`e $C=4\pi\epsilon_0a$ e il farad \`e la
capacit\`a di una sfera di raggio $9\times10^{11}$\,cm: \`e una unit\`a
decisamente troppo grossa, per cui di solito si usano i suoi sottomultipli: il
microfarad \ $(1\,\mu\mathrm F=10^{-6}\,\mathrm F)$, il nanofarad
\ $(1\,\mathrm{nF}=10^{-9}\,\mathrm F)$ ed il picofarad
\ $(1\,\mathrm{pF}=10^{-12}\,\mathrm F)$, che \`e praticamente uguale
all'unit\`a elettrostatica: $1\,\mathrm{cm}=1.1\,\mathrm{pF}$.

Siccome il numero di elettroni di conduzione in un conduttore non \`e
illimitato, potrebbe venirci il dubbio che non sia possibile caricare pi\`u che
tanto un conduttore: facciamo due conti all'ingrosso. Supponiamo di voler
portare una sferetta di 1\,cm di raggio al potenziale di $10^6$\,V, certamente
un valore spropositato. Siccome la capacit\`a \`e di 1\,pF, occorre la carica di
$10^{-6}$\,C, cio\`e occorre rimuovere dalla sferetta circa $10^{13}$ elettroni.
Quanti ce ne sono a disposizione? Contiamo solo quelli (di conduzione) vicino
alla superficie, perch\`e altrimenti la sferetta non si comporterebbe pi\`u come
un conduttore. Diciamo che la superficie ha uno spessore di $10^{-7}$\,cm,
allora, siccome in un metallo ci sono circa $10^{23}$ elettroni/cm$^3$, il
numero di elettroni a disposizione \`e
\ $4\pi\cdot10^{-7}\cdot10^{23}\simeq10^{17}$, quindi ci basta portarne via 1
ogni 10'000\,!

Consideriamo ora un sistema costituito da due conduttori: la relazione fra
cariche e potenziali comporta tre coefficienti di capacit\`a (o di potenziale),
tuttavia anche in questo caso \`e possibile, e soprattutto si rende necessario,
definire un'unica grandezza, detta capacit\`a del sistema. Supponiamo che i due
conduttori siano inizialmente scarichi, quindi, non essendoci altre cariche
intorno, sono entrambi a potenziale nullo. Se ora trasferiamo la carica $Q$ da
un conduttore all'altro, diciamo da $c_2$ a $c_1$, fra i due conduttori si
instaura una differenza di potenziale: la (2.9), che riscriviamo per il caso di
due conduttori, lega in generale le cariche ai potenziali
\begin{equation}
  \begin{aligned}
    \varphi_1&=P_{11}Q_1+P_{12}Q_2\\
    \varphi_2&=P_{21}Q_1+P_{22}Q_2
  \end{aligned}
  \tag{2.20}
\end{equation}
e, posto $V\equiv\varphi_1-\varphi_2$ (\`e la notazione pi\`u comunemente usata
in questo contesto),
\begin{equation}
  V=(P_{11}-P_{21})Q_1+(P_{12}-P_{22})Q_2
  \tag{2.21}
\end{equation}
per cui se $Q_1=Q$, \ $Q_2=-Q$
\begin{equation}
  V=(P_{11}+P_{22}-2P_{12})Q
  \tag{2.22}
\end{equation}
quindi, in queste condizioni, la differenza di potenziale fra i due conduttori
\`e proporzionale alla carica trasferita, per cui \`e possibile definire la
\textbf{capacit\`a del sistema}
\begin{equation}
  C\defeq\frac QV=(P_{11}+P_{22}-2P_{12})^{-1}
  \qquad(Q\equiv Q_1,\ V\equiv\varphi_1-\varphi_2)
  \tag{2.23}
\end{equation}
come il rapporto fra la carica trasferita e la differenza di potenziale.\\
Il sistema costituito dai due conduttori \`e detto \textbf{condensatore} (o
anche \textbf{capacitore}), in quanto -- quando $Q_2=-Q_1$ -- ``condensa il
campo elettrico'' nelle vicinanze del sistema (il campo lontano \`e un campo di
dipolo), cio\`e le linee di forza vanno da $c_1$ a $c_2$ senza disperdersi
all'infinito; i due conduttori sono detti le \textbf{armature} del
condensatore.\\
In termini dei coefficienti di capacit\`a
\begin{equation}
  C\,=\,\frac{C_{11}C_{22}-C_{12}^2}{C_{11}+C_{22}+2C_{12}}\,\cdot
  \tag{2.24}
\end{equation}
Se uno dei due conduttori \`e un conduttore cavo e l'altro, diciamo $c_1$, \`e
contenuto all'interno della cavit\`a, dalla discussione sulle gabbie di Faraday
sappiamo che la differenza di potenziale fra i due conduttori \`e indipendente
dalla carica $Q_2$ di $c_2$; quindi dalla (2.21) si ha
\begin{equation}
  P_{12}=P_{22}\quad\Rightarrow\quad C_{21}=-C_{11}
  \tag{2.25}
\end{equation}
ed in questo caso, che chiamiamo il caso del ``buon condensatore'', la
capacit\`a \`e data da
\begin{equation}
  C\,=\,(P_{11}-P_{21})^{-1}\,=\,C_{11}.
  \tag{2.26}
\end{equation}
Perch\`e ``buon condensatore''? Perch\`e la sua capacit\`a \`e ben definita e
resta invariata anche se il sistema si trova in presenza di altri conduttori o
cariche esterne, cosa che non succede per un condensatore le cui linee di forza
``escono allo scoperto'', e possono quindi venire modificate dalla presenza di
altri oggetti carichi.\\
Il condensatore costituito da una sfera al centro di una corona sferica
(``condensatore sferico''), che abbiamo discusso nell'ultimo paragrafo, \`e un
buon condensatore, e dalla (2.15) o dalla (2.16) la sua capacit\`a \`e data da
\begin{equation}
  C\,=\,\frac{ab}{b-a}
  \tag{2.27}
\end{equation}
da cui si vede che non dipende dal raggio esterno della corona.\\
Nelle applicazioni \`e importante disporre di condensatori di grande capacit\`a
e piccolo ingombro: accenneremo poi alla tecnica costruttiva; per ora notiamo
che dalla (2.27) risulta che la capacit\`a aumenta se diminuiamo la distanza fra
la sfera interna e la corona: se $\delta\equiv(b-a)\ll a$, dalla (2.27) si
ottiene
\begin{equation}
  C\,=\,\frac{a^2}{\delta}=\frac{S}{4\pi\delta}
  \tag{2.28}
\end{equation}
dove $S=4\pi a^2$ \`e la superficie del condensatore; la (2.28) \`e maggiore per
il fattore $a/\delta$ della capacit\`a di una sfera di raggio $a$ rispetto
all'infinito.\\
Consideriamo ora il condensatore costituito da due conduttori piani e paralleli
di area $S$, a distanza $\delta$ (``condensatore piano''). Supponiamo che
$\delta\ll\sqrt S$, in modo da poter approssimare il campo elettrico fra i due
piani come il campo elettrico (1.54) fra due distribuzioni piane infinite: si
dice che stiamo trascurando gli ``effetti di bordo'', un problema sul quale
torneremo pi\`u avanti. La differenza di potenziale fra le armature del
condensatore \`e data dalla (1.71), per cui dalla (2.23) si ha
\begin{equation}
  C=\frac{\sigma S}{4\pi\sigma\delta}=\frac{S}{4\pi\delta}
  \tag{2.29}
\end{equation}
cio\`e la stessa espressione del condensatore sferico.

La tecnologia per la costruzione di condensatori si evolve continuamente per
venire incontro alle sempre crescenti esigenze di miniaturizzazione dei
circuiti: l'idea di base \`e sostanzialmente quella di interporre fra due strati
conduttori sottili (anticamente fogli di alluminio), uno strato di ossido (o
altro isolante) quasi monomolecolare; in questo modo si ottiene un sistema che
ha una capacit\`a di almeno $10^{-2}\,\mathrm F/\mathrm m^2$. (Anticamente si
usava come strato isolante la carta che, per quanto fosse sottile, certamente
non era monomolecolare: questo ``sandwich'' veniva ripiegato e impacchettato
strettamente in un piccolo volume).

Nella figura 2.6 viene riportato il simbolo usato nei circuiti elettrici per
rappresentare un condensatore e il collegamento in serie e in parallelo di due
condensatori. Due condensatori sono collegati in serie quando una armatura
dell'uno \`e elettricamente collegata ad una armatura dell'altro: il
collegamento elettrico fra due elementi di un circuito viene rappresentata da
una linea sottile, in quanto normalmente \`e realizzata con fili conduttori. Il
collegamento in parallelo si realizza collegando ogni armatura di un
condensatore ad una armatura dell'altro.
\fig{fig2_6}{2.6}
Supponiamo che i due condensatori, di capacit\`a $C_1$ e $C_2$, siano ``buoni
condensatori'', per cui la presenza dell'uno non modifica la capacit\`a
dell'altro; quando essi vengono collegati in serie il sistema che ne risulta \`e
equivalente ad un condensatore di capacit\`a
\begin{equation}
  C=\Big(\frac{1}{C_1}+\frac{1}{C_2}\Big)^{-1}.
  \tag{2.30}
\end{equation}
Infatti, quando la carica $Q$ viene trasferita da una delle armature libere
all'altra, sulle armature collegate vengono indotte rispettivamente le cariche
$-Q$ e $Q$ (buoni condensatori!), e siccome
\[
  V=V_1+V_2=Q/C_1+Q/C_2
\]
ne segue la (2.30).
\fig{fig2_7}{2.7}
Se i due ``buoni'' condensatori vengono collegati in parallelo, la capacit\`a
del sistema \`e la somma delle capacit\`a:
\begin{equation}
  C=C_1+C_2
  \tag{2.31}
\end{equation}
in quanto la differenza di potenziale fra le armature di ciascun condensatore
\`e la stessa, e la carica \`e la somma delle cariche: $Q=Q_1+Q_2=VC_1+VC_2$, da
cui la (2.31).

Concludiamo questo paragrafo con l'espressione dell'energia di un condensatore
carico, cio\`e del lavoro che occorre fare per trasferire la carica $Q$ da una
armatura all'altra: dalla (1.92), usando la (2.23) si ottengono le seguenti
espressioni equivalenti
\begin{equation}
  U\equiv\frac12(Q_1\varphi_1+Q_2\varphi_2)
  =\frac12QV=\frac12CV^2=\frac12\,\frac{Q^2}{C}\quad\ (Q=Q_1=-Q_2).
  \tag{2.32}
\end{equation}
Nei vari casi particolari (p.es.\ condensatore sferico, condensatore piano) \`e
possibile verificare direttamente l'equivalenza delle (2.32) con la (1.94).

\subsection{Metodi di soluzione di problemi di elettrostatica dei conduttori}

L'idea che sta alla base di molti metodi di risoluzione di problemi di
elettrostatica dei conduttori \`e quella di sfruttare la conoscenza della
soluzione di un problema di elettrostatica, tipicamente di cariche fisse, per
risolverne un altro. Per esempio, sia data una distribuzione di cariche,
$\varphi(x,y,z)$ il potenziale prodotto da essa, e sia $S$ una superficie
equipotenziale $\varphi(x,y,z)=\varphi_0$. Dalla discussione sulle gabbie di
Faraday sappiamo che se $\varphi$ non \`e costante dentro $S$, \ $S$ \ deve
contenere delle cariche. Rimosse (solo) le cariche eventualmente contenute
dentro $S$, se \ $c$ \ \`e un conduttore la cui superficie esterna coincide con
$S$, \underline{all'esterno} del conduttore $\varphi(x,y,z)$ risolve il problema
di Dirichlet con il conduttore a potenziale $\varphi_0$: infatti
$\varphi(x,y,z)$ soddisfa l'equazione di Poisson nella regione esterna al
conduttore, e la corretta condizione al contorno su $S$. Per l'unicit\`a della
soluzione, \`e la soluzione. Inoltre, la densit\`a di carica sulla superficie
del conduttore \`e data dalla (2.2) e quindi, essendo noto il potenziale, \`e
calcolabile; il campo elettrico su $S$, e quindi anche il flusso di $\vec E$
uscente da $S$, \`e lo stesso prima e dopo, quindi per il teorema di Gauss la
carica totale sul conduttore \`e uguale alla carica totale della distribuzione
che era contenuta dentro $S$ e che abbiamo rimosso.\\
Il metodo delle cariche immagine utilizza questo procedimento all'inverso: sia
dato un sistema di conduttori $c_i$ di superfici $S_i$, ed una distribuzione di
cariche nella regione $\mathcal R$ delimitata dalle superfici $S_i$ dei
conduttori e dalla ``massa'' $S_0$. Il problema \`e quello di determinare il
potenziale in $\mathcal R$, date le condizioni al contorno sui conduttori,
cio\`e per ciascuno di essi il potenziale $\varphi_i$ o la carica $Q_i$. Il
problema \`e risolto se \`e possibile trovare una distribuzione di cariche fisse
\underline{esterna} ad $\mathcal R$, cio\`e nelle regioni interne alle $S_i$
(\textbf{cariche immagine}), che, insieme alla distribuzione di cariche fisse
contenuta in $\mathcal R$, genera \underline{in assenza} dei conduttori $c_i$ un
potenziale $\varphi(x,y,z)$ che \`e costante su ogni $S_i$, e soddisfa su di
esse le condizioni al contorno assegnate: $\varphi(x,y,z)=\varphi_i$, se \`e
assegnato il potenziale, altrimenti $\Phi_{S_i}(\vec E)=4\pi Q_i$, se \`e
assegnata la carica. In tal caso $\varphi(x,y,z)$ \`e la soluzione in
$\mathcal R$ del problema con i conduttori.

\paragraph{Esempi}
\begin{enumerate}
  \item Consideriamo un dipolo: due cariche $q$ e $-q$ nelle posizioni $x=y=0$,
    $z=\pm a$.\\
    Il potenziale \`e
    \begin{equation}
      \varphi(x,y,z)=q\Big(\frac{1}{\sqrt{x^2+y^2+(z-a)^2}}
      -\frac{1}{\sqrt{x^2+y^2+(z+a)^2}}\Big)
      \tag{2.33}
    \end{equation}
    e assume tutti i valori da $-\infty$ (avvicinandosi a $-q$) a $+\infty$
    (avvicinandosi a $q$); il piano $z=0$ \`e la superficie equipotenziale con
    $\varphi=0$. Le superfici equipotenziali con $\varphi<0$ sono superfici
    chiuse nella regione $z<0$ contenenti la carica negativa (appiattite dalla
    parte del piano $z=0$, allungate nella direzione opposta), viceversa per
    quelle su cui $\varphi>0$. Se sostituiamo (per esempio) la carica negativa
    con un conduttore $c$ avente come superficie (esterna) una superficie
    equipotenziale $S$ che la contiene, se su $S$ \ $\varphi(x,y,z)=\varphi_0$
    \ la (2.33) risolve (esternamente al conduttore) il problema dell'induzione
    esercitata dalla carica $q$ sul conduttore $c$ il cui potenziale sia
    $\varphi_0$ (problema di Dirichlet). La carica totale sul conduttore $c$ (a
    potenziale $\varphi_0$) \`e $-q$.
    \fig{fig2_8}{2.8}
    Ovviamente la (2.33) risolve anche il problema (di Neumann) in cui sia
    assegnata la carica $-q$ del conduttore $c$. In entrambi i casi la carica
    $-q$ (che non c'\`e pi\`u) \`e la carica immagine del problema ed \`e
    situata esternamente alla regione che ci interessa.\\
    In particolare, e questo \`e il caso pi\`u interessante, il conduttore pu\`o
    essere una lastra piana infinita la cui superficie ``esterna'' coincide con
    il piano $z=0$. In questo caso la (2.33) risolve, nella regione $z\ge0$, il
    problema elettrostatico (che abbiamo gi\`a in parte discusso) di una carica
    puntiforme $q$ posta a distanza $a$ dalla superficie di un conduttore piano
    infinito, messo a terra ($\varphi_0=0$). Anche in questo caso la carica $-q$
    \`e la carica immagine del problema ed \`e situata esternamente alla regione
    che ci interessa ($z\ge0$).\\
    La carica indotta sulla superficie del conduttore si calcola facilmente:
    \`e, per la (2.2), proporzionale al campo elettrico associato al potenziale
    (2.33), sul piano $z=0$
    \begin{equation}
      \sigma(x,y)=\frac{1}{4\pi}\,E_z(x,y,0)
      =-\frac{1}{4\pi}\,\frac{2qa}{(x^2+y^2+a^2)^{3/2}}\,\cdot
      \tag{2.34}
    \end{equation}
    La carica totale sulla superficie $z=0$ del conduttore si pu\`o ottenere
    dalla (2.34) per integrazione, oppure mediante il teorema di Gauss applicato
    alla superficie infinita costituita dal piano $z=-\epsilon$ chiuso dalla
    semisfera all'infinito nella regione della carica $+q$ (pi\`u correttamente,
    figura 2.9, si considera il flusso attraverso la semisfera di raggio $R$ e
    centro nell'origine, completata dal disco equatoriale di raggio $R$ appena
    dentro il conduttore, e poi si prende il limite per $R\to\infty$): detto
    flusso \`e nullo perch\`e \`e nullo il campo elettrico dentro il conduttore,
    ed \`e nullo il flusso attraverso la semisfera all'infinito perch\`e il
    campo \`e un campo di dipolo, e quindi di ordine $O(r^{-3})$ all'infinito,
    mentre la superficie \`e di ordine $O(r^2)$. Quindi la carica totale sulla
    superficie del conduttore \`e $-q$.
    \fig{fig2_9}{2.9}
    Vale la pena di osservare che il momento di dipolo della distribuzione di
    carica ($q$ e $\sigma$) \`e parallelo all'asse $z$ (grazie alla simmetria
    delle cariche) e vale $p=qa$ (basta
    prendere come polo l'origine e calcolare solo $p_z$), mentre il potenziale
    all'infinito \`e un potenziale di dipolo con $p=2qa$: la ragione della
    discrepanza sta nel fatto che la distribuzione delle sorgenti ($\sigma$) si
    estende su una regione infinita dello spazio, per cui il procedimento
    segu\`ito per ricavare lo sviluppo del potenziale a grandi distanze non
    pu\`o essere applicato.\\
    Sulla carica $q$ si esercita una forza dovuta al campo elettrico generato
    dalla carica sulla superficie del conduttore, e quindi uguale alla forza
    esercitata dalla carica immagine: quindi il conduttore attira la carica con
    una forza pari a
    \begin{equation}
      F=\frac{q^2}{4a^2}\,\cdot
      \tag{2.35}
    \end{equation}
    L'energia elettrostatica del sistema, viste le (1.89) e (1.92), pu\`o essere
    calcolata come
    \begin{equation}
      U=\frac12\Big(q\varphi^{(1)}(0,0,a)
      +\int\sigma(x,y)\varphi(x,y,0)dx\,dy\Big)=-\frac{q^2}{4a}
      \tag{2.36}
    \end{equation}
    dato che $\varphi^{(1)}$, come nella (1.89), \`e il potenziale dovuto solo
    alla distribuzione $\sigma(x,y)$\,: \ $\varphi^{(1)}(0,0,a)=-q/2a$, \ e
    \ $\varphi(x,y,0)=0$.\\
    Possiamo anche calcolare $U$ in modo pi\`u diretto come il lavoro necessario
    per portare la carica $q$ da $z=\infty$ a $z=a$ \ $(x=y=0)$: siccome quando
    la carica \`e a distanza $z$ dal piano su di essa, per la (2.35), si
    esercita la forza $F_z(z)=-q^2/4z^2$, si ha
    \begin{equation}
      U=-\int_\infty^aF_z(z)\,dz
      =\int_\infty^a\frac{q^2}{4z^2}\,dz=-\frac{q^2}{4a}\,\cdot
      \tag{2.36'}
    \end{equation}
    Si noti che $U$ \`e la met\`a dell'energia $-q^2/2a$ del dipolo costituito
    da $q$ e da $-q$ a distanza $2a$: infatti il lavoro per portare due cariche
    (reali) $q$ e $-q$ a distanza $2a$, mantenendole in posizione simmetrica
    rispetto al piano $z=0$, \`e il doppio del lavoro su ciascuna carica; nel
    caso che stiamo considerando una delle due cariche non c'\`e, e quindi il
    lavoro che occorre fare \`e la met\`a (oppure: quando viene spostata la
    carica reale, la carica immagine si sposta da sola!).
  \item Campo di una carica puntiforme $q$ a distanza $l$ dal centro $O$ di una
    sfera conduttrice di raggio $R<l$: \`e un altro problema classico del metodo
    delle cariche immagine.\\
    Cominciamo con il caso in cui la sfera \`e a potenziale zero (sfera a
    massa). Ovviamente ci interessa il potenziale all'esterno della sfera.
    Possiamo sperare di poter sostituire la sfera con una carica puntiforme $q'$
    a distanza $l'$ da $O$, con $l'<R$, in modo tale che il potenziale prodotto
    da $q$ e da $q'$ sia nullo su quella che era la superficie della sfera
    conduttrice? Proviamo: dette $r$ ed $r'$ le distanze di un punto della sfera
    da $q$ e $q'$, deve essere \ $q/r+q'/r'=0$, cio\`e \ $r/r'=-q/q'$; ma questa
    \`e proprio la definizione di cerchio di Apollonio, nel piano, e di sfera di
    Apollonio, nello spazio. Il cerchio (o la sfera) di Apollonio ha il centro
    sul prolungamento della retta congiungente $q$ e $q'$ nel punto a distanza
    $l$ da $q$ ed $l'$ da $q'$ tali
    \fig{fig2_10}{2.10}
    che $l/l'=(q/q')^2$, e raggio $R=\sqrt{ll'}$. Quindi, i dati essendo $R$\,,
    $q$ ed $l$, una carica $q'=-q\,R/l$ a distanza $l'=R^2/l$ dal centro della
    sfera, \`e la carica immagine del problema.\\
    Il potenziale in un punto $P$ esterno al conduttore \`e
    \begin{equation}
      \varphi(x,y,z)=\frac qr-\frac{qR}{lr'}
      \tag{2.37}
    \end{equation}
    dove $r$ ed $r'$, indicate nella figura 2.10, sono le distanze di $P$ da $q$
    e dalla posizione della carica immagine $q'$. La carica totale indotta sulla
    sfera \`e $q'=-q\,R/l$. Si noti che se $q$ si allontana, $q'$ si avvicina al
    centro.\\
    Supponiamo ora che la sfera, anzich\`e a potenziale zero, sia una sfera
    scarica e isolata, cio\`e la sua carica totale \`e nulla: se aggiungiamo al
    sistema un'altra carica immagine $-q'$ nel centro della sfera, la superficie
    della sfera resta equipotenziale (a potenziale $\varphi_0=-q'/R=q/l$), e la
    carica totale sulla sfera \`e nulla. In questo caso il potenziale fuori dal
    conduttore \`e
    \begin{equation}
      \varphi(x,y,z)=\frac qr-\frac{qR}{lr'}+\frac{qR}{lr_0}
      \tag{2.37}
    \end{equation}
    dove $r_0$ \`e la distanza del punto $P$ dal centro della sfera. In
    generale, una opportuna carica nel centro permette di soddisfare qualsiasi
    condizione al contorno sulla sfera.
  \item Sfera conduttrice in un campo elettrico uniforme. Vogliamo stabilire
    come un campo uniforme $\vec E_0$, per esempio generato da due piani
    paralleli con densit\`a di carica (fissa) $\sigma_0$ e $-\sigma_0$, viene
    modificato dalla presenza di una sfera conduttrice di raggio $R$ scarica e
    isolata, e determinare la densit\`a di carica indotta sulla sua superficie.
    Si pu\`o procedere in due modi:

    \textit{i}) \ prendiamo l'asse $z$ parallelo e concorde con $\vec E_0$; il
    potenziale \`e \ $\varphi_0(x,y,z)=-E_0z$, e sulla superficie della sfera
    vale $-E_0R\cos\theta$. Le eventuali cariche immagine (carica totale nulla)
    devono produrre sulla superficie della sfera un potenziale uguale e opposto
    (tutt'al pi\`u a meno di una costante): il potenziale di un dipolo
    puntiforme \`e, ad ogni distanza,
    \[
      \frac{\vec p\cdot\vec r}{r^3}=\frac{p\cos\theta}{r^2}
    \]
    per cui, se prendiamo come cariche immagine un dipolo puntiforme con
    $\vec p=R^3\vec E_0$ posto nel centro della sfera, il potenziale
    \begin{equation}
      \varphi(x,y,z)=-E_0z+\frac{\vec p\cdot\vec r}{r^3}
      =-E_0z\Big(1-\frac{R^3}{r^3}\Big)
      =-E_0r\cos\theta\Big(1-\frac{R^3}{r^3}\Big)
      \tag{2.38}
    \end{equation}
    \`e nullo sulla superficie della sfera, e per $r\ge R$ risolve il problema.

    \textit{ii}) \ Possiamo pensare il campo uniforme $\vec E_0$ come prodotto
    da due cariche $q$ e $-q$ nelle posizioni $x=y=0$, $z=\mp z_0$ nel limite
    $z_0\to\infty$, $q\to\infty$: infatti \big(il calcolo \`e simile a quello
    che ha portato alla (1.75)\big)
    \[
      \varphi_0(x,y,z)\equiv q\Big(\frac{1}{\sqrt{x^2+y^2+(z+z_0)^2}}
      -\frac{1}{\sqrt{x^2+y^2+(z-z_0)^2}}\Big)
    \]
    \begin{equation}
      =\frac{q}{z_0}\Big(\frac{1}{\sqrt{1+(r/z_0)^2+2z/z_0}}
      -\frac{1}{\sqrt{1+(r/z_0)^2-2z/z_0}}\Big)\to-\frac{2qz}{z_0^2}
      \tag{2.39}
    \end{equation}
    per cui basta che $2q/z_0^2\to E_0$ quando $z_0\to\infty$. Le cariche
    immagine (rispetto alla sfera conduttrice) delle due cariche $q$ e $-q$,
    sono le cariche $-qR/z_0$ e $+qR/z_0$ nei punti $\mp R^2/z_0$ ed hanno
    momento di dipolo $p=2qR^3/z_0^2$ concorde con l'asse $z$. Quando
    $z_0\to\infty$ le due cariche immagine costituiscono un dipolo puntiforme
    con $\vec p=R^3\vec E_0$ posto nel centro della sfera, come prima.

    Per determinare la carica indotta sulla superficie della sfera occorre
    calcolare il campo elettrico sulla sfera: siccome sappiamo che $\vec E$ \`e
    ortogonale alla superficie, \`e sufficiente che calcoliamo la derivata di
    $\varphi$ rispetto ad $r$: per $r$ costante, uguale ad $R$, $\varphi$ non
    varia. Dalla (2.38) si ha
    \begin{equation}
      \begin{aligned}
        4\pi\sigma=E_n&=-\frac{\partial\varphi}{\partial r}=3E_0\cos\theta
        \quad\Rightarrow\\
        \Rightarrow\sigma&=\frac{3}{4\pi}\,E_0\cos\theta\equiv\sigma_0\cos\theta
        \qquad\quad\sigma_0\equiv\frac{3E_0}{4\pi}
      \end{aligned}
      \tag{2.40}
    \end{equation}
    quindi il campo elettrico $\vec E_0$ induce sulla sfera la distribuzione
    (2.40) che ha momento di dipolo $\vec p=R^3\vec E_0$. Siccome il campo
    dentro la sfera conduttrice \`e nullo, segue che la distribuzione (2.40)
    produce per $r<R$ un campo elettrico uniforme pari a $-\vec E_0$ quindi, in
    generale, una distribuzione superficiale $\sigma_0\cos\theta$ sulla
    superficie di una sfera, produce dentro la sfera il campo
    \ $\vec E=-(4\pi/3)\sigma_0\,\hat z$.
\end{enumerate}

\subsection{Effetti di bordo}

Abbiamo calcolato il campo elettrico di distribuzioni di carica aventi come
supporto cilindri infiniti, fili rettilinei, piani (conduttori e non) infiniti,
ecc., ben sapendo che tali distribuzioni non corrispondono a situazioni reali:
sembra quasi che l'unica ragione per cui ce ne siamo occupati sia quella che in
questi casi campo elettrico e potenziale, grazie alle simmetrie di questi
sistemi, sono facilmente calcolabili. In un certo senso questo \`e vero, per\`o
non \`e l'unica ragione: lo vedremo esplicitamente su alcuni esempi, ma \`e
abbastanza chiaro a priori che i campi calcolati per queste distribuzioni
possono costituire una buona approssimazione dei campi di distribuzioni finite,
in opportune regioni di spazio di dimensioni piccole rispetto alle dimensioni
delle distribuzioni. Per esempio, abbiamo visto alla fine del capitolo 1.10 che
una distribuzione lineare di lunghezza finita ha la stessa singolarit\`a di una
distribuzione lineare su un filo rettilineo infinito, quindi a piccole distanze
le due distribuzioni diventano indistinguibili: la ragione di ci\`o sta nella
legge $1/r^2$, per cui i contributi delle regioni lontane di una distribuzione
diventano sempre meno importanti. Per lo stesso motivo ci aspettiamo che una
distribuzione costante $\sigma$ sulla superficie di un disco di raggio $R$
generi un campo elettrico che in prossimit\`a del centro del disco non si
discosti molto da quello di una identica distribuzione su un piano infinito.
Supponiamo il disco nel piano $z=0$ con il centro nell'origine, il campo
elettrico pu\`o essere calcolato facilmente sull'asse (asse $z$), ed \`e
\begin{equation}
  E_x^0=E_y^0
  =0,\ \ E_z^0=2\pi\sigma\Big(\frac{z}{|z|}-\frac{z}{\sqrt{z^2+R^2}}\Big)
  \tag{2.41}
\end{equation}
Per $x=y=0$, \ $z=0^\pm$ \ $E_z=\pm2\pi\sigma$, come per una distribuzione piana
infinita; in un intorno dell'origine, per $z>0$ si ha
\[
  E_z=2\pi\sigma\Big(1-\frac zR+O\big(z^3/R^3\big)\Big)
\]
che conferma che l'approssimazione del disco con un piano infinito \`e
accettabile vicino al centro finch\`e $z\ll R$. Pi\`u ci si allontana dal
centro, sia in direzione dell'asse $z$, sia verso il bordo, pi\`u il campo si
discosta da quello di un piano infinito: questo scostamento \`e ci\`o che
normalmente viene indicato come ``effetto di bordo''. Abbiamo trascurato gli
effetti di bordo quando abbiamo calcolato la capacit\`a di un condensatore
piano: in questo caso gli effetti di bordo non sono dovuti soltanto alla
dimensione finita della distribuzione, ma anche al fatto che, siccome le
armature sono dei conduttori, la carica non \`e distribuita uniformemente, come
in un piano infinito, ma si ha un accumulo verso i bordi. Il campo elettrico non
\`e completamente confinato fra le armature, ma \`e diverso da zero anche
all'esterno e in particolare, oltre alle linee di forza che uniscono le facce
interne delle armature, ci sono linee di forza che vanno dall'esterno di
un'armatura all'esterno dell'altra, e queste si hanno principalmente verso i
bordi delle armature.\\
La necessit\`a di un campo elettrico anche all'esterno del condensatore \`e
dettata dal fatto che il campo elettrico \`e conservativo: se $\vec E$ fosse
confinato in una regione limitata di spazio, la circuitazione di $\vec E$ lungo
un circuito del tipo di quello rappresentato in figura 2.11 (il lato interno
lungo una linea di forza, uno esterno dove ipotizziamo che il campo sia nullo, e
gli altri due ortogonali ad $\vec E$ dove $\vec E\neq0$) sarebbe diversa da
zero, in quanto riceverebbe contributo solo dal lato interno (in effetti il
campo a grande distanza dal condensatore \`e un campo di dipolo, con momento di
dipolo pari alla carica su un'armatura per la distanza fra le armature).
\fig{fig2_11}{2.11}
\`E chiaro quindi che trascurare gli effetti di bordo pu\`o portare a
conclusioni in contrasto con la conservazione dell'energia: \`e quello che
accade se un condensatore tipo quello di figura 2.11 viene attraversato da una
particella carica: se il campo elettrico fosse diverso da zero solo dentro il
condensatore, e quindi non conservativo, la particella verrebbe accelerata (in
direzione del campo) attraversando il condensatore, ed uscirebbe con energia
cinetica diversa da quella iniziale; si potrebbe poi riportare la particella
nella posizione iniziale lungo un percorso tutto esterno al condensatore dove il
campo elettrico \`e (per ipotesi) nullo. Nella realt\`a, invece, le linee di
forza vicino ai bordi, e soprattutto all'esterno del condensatore, provvedono a
rallentare la particella in modo che all'inizio e alla fine abbia la stessa
energia cinetica.\\
Questi effetti di bordo portano a correzioni alla capacit\`a di un condensatore
piano che, nel caso di armature circolari di raggio $R$, risulta maggiore di
quella data dalla (2.29) per circa il 17\% per $\delta/R=0.1$, e di circa il 2\%
se $\delta/R=0.01$.
